\newcommand{\titulus}{\nomenFesti{S. Athanasii, Episcopi \& Ecclesiæ Doctoris.}
\dies{Die 2. Maii.}}
\newcommand{\oratio}{\pars{Oratio.}

\noindent Omnípotens sempitérne Deus, qui beátum Athanásium, epíscopum, divinitátis Fílii tui propugnatórem exímium suscitásti, concéde propítius, ut, eius doctrína et protectióne gaudéntes, in tui cognitióne et amóre sine intermissióne crescámus.

\pars{Pro pace in universo mundo.} \scriptura{Sir. 50, 25; 2 Esdr. 4, 20; \textbf{H416}}

\vspace{-4mm}

\antiphona{II D}{temporalia/ant-dapacemdomine.gtex}

\vfill

\noindent Deus, a quo sancta desidéria, recta consília et iusta sunt ópera: da servis tuis illam, quam mundus dare non potest, pacem; ut et corda nostra mandátis tuis dédita, et hóstium subláta formídine, témpora sint tua protectióne tranquílla.

\noindent Per Dóminum nostrum Iesum Christum, Fílium tuum, qui tecum vivit et regnat in unitáte Spíritus Sancti, Deus, per ómnia sǽcula sæculórum.

\noindent \Rbardot{} Amen.}
\newcommand{\invitatorium}{\pars{Invitatorium.}

\vspace{-4mm}

\antiphona{VI*}{temporalia/inv-fontemsapientiae-tp.gtex}}
\newcommand{\hymnusmatutinum}{\pars{Hymnus}

\cuminitiali{III}{temporalia/hym-ChristePastorum.gtex}}
\newcommand{\lectioi}{\pars{Lectio I.} \scriptura{Ap. 18, 1-8}

\noindent De libro Apocalýpsis beáti Ioánnis apóstoli.

\noindent Ego Ioánnes vidi álium ángelum descendéntem de cælo, habéntem potestátem magnam, et terra illumináta est a claritáte eius.

\noindent Et clamávit in forti voce dicens: «Cécidit, cécidit Bábylon magna et facta est habitátio dæmoniórum et custódia omnis spíritus immúndi et custódia omnis béstiæ immúndæ et odíbilis, quia de vino iræ fornicatiónis eius bibérunt omnes gentes, et reges terræ cum illa fornicáti sunt, et mercatóres terræ de virtúte deliciárum eius dívites facti sunt!».

\noindent Et audívi áliam vocem de cælo dicéntem: «Exíte de illa, pópulus meus, ut ne compartícipes sitis peccatórum eius et de plagis eius non accipiátis, quóniam pervenérunt peccáta eius usque ad cælum, et recordátus est Deus iniquitátum eius. Réddite illi, sicut et ipsa réddidit, et duplicáte duplícia secúndum ópera eius; in póculo, quo míscuit, miscéte illi duplum. Quantum glorificávit se et in delíciis fuit, tantum date illi torméntum et luctum. Quia in corde suo dicit: “Sédeo regína et vídua non sum et luctum non vidébo”, ídeo in una die vénient plagæ eius, mors et luctus et fames, et igne comburétur, quia fortis est Dóminus Deus, qui iudicávit illam».}
\newcommand{\responsoriumi}{\pars{Responsorium 1.} \scriptura{\Rbardot{} Cf. 1 Chr. 13, 8 \Vbardot{} Ps. 98, 6; \textbf{H248}}

\vspace{-5mm}

\responsorium{VI}{temporalia/resp-decantabatpopulus-CROCHU.gtex}{}}
\newcommand{\lectioii}{\pars{Lectio II.} \scriptura{Ap. 18, 9-20}

\noindent Et flebunt et plangent se super illam reges terræ, qui cum illa fornicáti sunt et in delíciis vixérunt, cum víderint fumum incéndii eius, longe stantes propter timórem tormentórum eius, dicéntes: «Væ, væ, cívitas illa magna, Bábylon, cívitas illa fortis, quóniam una hora venit iudícium tuum!».

\noindent Et negotiatóres terræ flent et lugent super illam, quóniam mercem eórum nemo emit ámplius: mercem auri et argénti et lápidis pretiósi et margaritárum et byssi et púrpuræ et sérici et cocci, et omne lignum thýinum et ómnia vasa éboris et ómnia vasa de ligno pretiosíssimo et æraménto et ferro et mármore, et cinnamómum et amómum et odoraménta et unguénta et tus, et vinum et óleum et símilam et tríticum, et iuménta et oves et equórum et rædárum, et mancipiórum et ánimas hóminum.

\noindent Et fructus tui, desidérium ánimæ, discessérunt a te, et ómnia pínguia et clara periérunt a te, et ámplius illa iam non invénient.

\noindent Mercatóres horum, qui dívites facti sunt ab ea, longe stabunt propter timórem tormentórum eius flentes ac lugéntes, dicéntes: «Væ, væ, cívitas illa magna, quæ amícta erat býssino et púrpura et cocco, et deauráta auro et lápide pretióso et margaríta, quóniam una hora desolátæ sunt tantæ divítiæ!».

\noindent Et omnis gubernátor et omnis, qui in locum návigat, et nautæ et, quotquot mária operántur, longe stetérunt et clamábant, vidéntes fumum incéndii eius, dicéntes: «Quæ símilis civitáti huic magnæ?».

\noindent Et misérunt púlverem super cápita sua et clamábant flentes et lugéntes, dicéntes: «Væ, væ, cívitas illa magna, in qua dívites facti sunt omnes, qui habent naves in mari, de ópibus eius, quóniam una hora desoláta est! Exsúlta super eam, cælum, et sancti et apóstoli et prophétæ, quóniam iudicávit Deus iudícium vestrum de illa!».}
\newcommand{\responsoriumii}{\pars{Responsorium 2.} \scriptura{\Rbardot{} Ap. 21, 2.18 \Vbardot{} Ps. 44, 14; \textbf{H248}}

\vspace{-5mm}

\responsorium{I}{temporalia/resp-vidijerusalemdescendentem-CROCHU.gtex}{}}
\newcommand{\lectioiii}{\pars{Lectio III.} \scriptura{Oratio de incarnatione Verbi, 8-9: PG 25, 110-111}

\noindent Ex Oratiónibus sancti Athanásii epíscopi.

\noindent {\color{gray} Incorpóreum et corruptiónis ac matériæ expers, Dei Verbum in nostram regiónem advénit, quamquam ántea non procul áberat: nulla enim pars mundi illo umquam vácua fuit, sed una cum suo Patre exsístens ómnia ubíque implébat.

\noindent Venit vero sua erga nos benignitáte, et quátenus sese palam nobis exhíbuit. Nostri ipse géneris et infirmitátis misértus, nostráque motus corruptióne, nec mortem in nobis dominári ferens, ne scílicet quod factum fúerat períret, suíque Patris in hómine formándo opus vanum fíeret, sibi ipsi corpus accépit, idque nostri non dissímile; nec enim esse tantum in córpore, vel solum apparére vóluit.

\noindent Nam, si tantum apparére voluísset, potuísset sane áliud præstántius corpus assúmere; verum nostrum corpus accépit.}

\noindent Sibi in Vírgine templum, corpus scílicet, exstrúxit, illúdque tamquam instruméntum sibi próprium fecit, in quo se notum fáceret et habitáret: sic ergo símili córpore e nostris accépto, quia omnes mortis corruptióni subiécti erant, illud, pro ómnibus morti tráditum, Patri summa cum benignitáte óbtulit; cum, ut ómnibus in ipso moriéntibus lex de corruptióne contra hómines lata solverétur, útpote quæ in domínico córpore suam vim complevísset, nec proínde locum ámplius advérsus símiles hómines habéret; tum, ut hómines, in corruptiónem revérsos, íterum incorrúptos rédderet, atque a morte ad vitam revocáret, córpore quod sibi assúmpserat et resurrectiónis grátia mortem ab eis, non secus ac stípula ab igne consúmitur, pénitus ámovens.

\noindent Idcírco corpus quod mori posset sibi assúmpsit, ut illud, Verbi ómnium prǽsidis factum párticeps, morti pro ómnibus satis esset, atque, propter Verbum in se hábitans, incorrúptum permanéret, ac dénique in pósterum corrúptio resurrectiónis grátia ab ómnibus discéderet. Hinc corpus, quod sibi ipse accépit, velut hóstiam et víctimam omni puram mácula morti offeréndo, mortem statim ab ómnibus simílibus, suo pro áliis obláto, propulsávit.

\noindent Sic enim Dei Verbum, supérius ómnibus exsístens, suum templum et corpóreum instruméntum pro ómnibus, uti dicébat, devovéndo et offeréndo, id quod debebátur in morte solvit, atque ita, ómnibus per símile corpus coniúnctus, incorrúptus Dei Fílius omnes resurrectiónis promissióne incorrúptos iure et mérito réddidit. Ipsa síquidem mortis corrúptio nullam iam vim advérsus hómines habet ob Verbum, quod in illis per unum corpus hábitat.}
\newcommand{\responsoriumiii}{\pars{Responsorium 3.} \scriptura{\Rbardot{} Ps. 63, 11 \Vbardot{} ibid., 10; \textbf{H253}}

\vspace{-5mm}

\responsorium{V}{temporalia/resp-laetabituriustus-CROCHU-cumdox.gtex}{}}
\newcommand{\hymnuslaudes}{\pars{Hymnus}

\cuminitiali{VIII}{temporalia/hym-InclitusPastor-doctorum.gtex}}
\newcommand{\benedictus}{\pars{Canticum Zachariæ.} \scriptura{Mt. 7, 15.16; \textbf{H429}}

\vspace{-6mm}

{
\grechangedim{interwordspacetext}{0.18 cm plus 0.15 cm minus 0.05 cm}{scalable}%
\antiphona{VIII G}{temporalia/ant-attenditeafalsisprophetis.gtex}
\grechangedim{interwordspacetext}{0.22 cm plus 0.15 cm minus 0.05 cm}{scalable}%
}

\vspace{-4mm}

\scriptura{Lc. 1, 68-79}

\vspace{-2mm}

\cantusSineNeumas
\initiumpsalmi{temporalia/benedictus-initium-viiisoll-G-auto.gtex}

\vspace{-1.5mm}

\input{temporalia/benedictus-viiisoll-G.tex} \Abardot{}
}
\include{hebdomadaiv}
\include{sabbato}
